\documentclass{article}
\usepackage[margin=1in]{geometry}
\usepackage{siunitx}
\usepackage{amsmath}
\sisetup{per-mode=symbol}
\title{Density Calculations Worksheet}
\author{}
\date{}
\begin{document}
\maketitle

\section*{Questions}

\begin{enumerate}
\item A lump of steel has a mass of \SI{7600}{\kilogram} and a volume of \SI{1}{\metre\cubed}. What is the density of steel?

\item A container of water holds a volume of \SI{2}{\metre\cubed} of water; it weighs \SI{2000}{\kilogram}. What is the density of water?

\item A lump of copper has a mass of \SI{18400}{\kilogram} and a volume of \SI{2.3}{\metre\cubed}. What is the density of copper?

\item Brass has a density of \SI{7500}{\kilogram\per\metre\cubed}. What is the mass if you have \SI{2.1}{\metre\cubed} of it?

\item Lead is very dense with a density of \SI{13000}{\kilogram\per\metre\cubed}. If you have a sample with a mass of \SI{2000}{\kilogram}, what is the volume of this sample?

\item Gold has a mass of \SI{45}{\kilogram}.  It has a density of \SI{19300}{\kilogram\per\metre\cubed}.  What is the volume of this gold?

\item A \emph{cylinder} of brass \(\bigl(\text{density} = \SI{7500}{\kilogram\per\metre\cubed}\bigr)\) has a radius of \SI{0.01}{\metre} and a length of \SI{0.05}{\metre}. What is its mass?

\item A \emph{cube} of steel has a density of \SI{7860}{\kilogram\per\metre\cubed}. If each side is \SI{0.10}{\metre} long, what is the total mass of the steel cube?

\item A \emph{cuboid} made of glass has the following dimensions: width \SI{20}{\centi\metre}, height \SI{20}{\centi\metre}, length \SI{8}{\metre}.  Calculate the density of this type of glass if the mass of this cuboid is \SI{800}{\kilogram}.

\item A steel \emph{ball bearing} has a radius of \SI{0.5}{\centi\metre}. Steel has a density of \SI{7860}{\kilogram\per\metre\cubed}. Calculate the mass of the ball bearing.

\end{enumerate}

\subsection*{Challenge Questions}

\begin{enumerate}
\setcounter{enumi}{10}
\item A \emph{concrete triangular prism} is \SI{5}{\metre} long with a triangular cross‐section of base \SI{6}{\metre}. Its mass is \SI{288000}{\kilogram}.  Calculate the \emph{height} of the prism.\par
\textit{(Density of concrete: \SI{2400}{\kilogram\per\metre\cubed}.)}

\item A solid beryllium sphere has a mass of \SI{3966}{\kilogram}.  Given that beryllium has a density of \SI{1850}{\kilogram\per\metre\cubed}, calculate the radius of the sphere used in the particle accelerator.

\item A \emph{hollow} spherical geode has an outer radius of \SI{4}{\centi\metre} and a crust \SI{0.5}{\centi\metre} thick made of limestone \(\bigl(\text{density } \SI{1760}{\kilogram\per\metre\cubed}\bigr)\). Calculate the mass of just the limestone crust.
\end{enumerate}

\vspace{1em}
\noindent\textbf{Useful formulae}\\[0.3em]
\begin{tabular}{ll}
Volume of a cube: & $V = a^{3}$\\
Volume of a cuboid: & $V = l w h$\\
Volume of a cylinder: & $V = \pi r^{2} h$\\
Volume of a triangular prism: & $V = \tfrac12 b h \, l$\\
Volume of a sphere: & $V = \tfrac43 \pi r^{3}$\\
Density relationship: & $\rho = \dfrac{m}{V}$
\end{tabular}

\newpage
\section*{Answer Key}

\begin{enumerate}
\item $\rho = \dfrac{7600}{1} = \SI{7600}{\kilogram\per\metre\cubed}$

\item $\rho = \dfrac{2000}{2} = \SI{1000}{\kilogram\per\metre\cubed}$

\item $\rho = \dfrac{18400}{2.3} = \SI{8000}{\kilogram\per\metre\cubed}$

\item $m = \rho V = 7500 \times 2.1 = \SI{15750}{\kilogram}$

\item $V = \dfrac{m}{\rho} = \dfrac{2000}{13000} = \SI{0.154}{\metre\cubed}$

\item $V = \dfrac{45}{19300} = \SI{0.00233}{\metre\cubed}$

\item $V = \pi r^{2} h = \pi(0.01)^{2}(0.05) = \SI{1.58e-5}{\metre\cubed}$\\
      $m = \rho V = 7500 \times 1.58\times10^{-5} = \SI{0.118}{\kilogram}$

\item $V = a^{3} = (0.10)^{3} = \SI{1.00e-3}{\metre\cubed}$\\
      $m = \rho V = 7860 \times 1.00\times10^{-3} = \SI{7.86}{\kilogram}$

\item Convert the widths: \SI{20}{\centi\metre} = \SI{0.20}{\metre}.\\
      $V = l w h = 8 \times 0.20 \times 0.20 = \SI{0.32}{\metre\cubed}$\\
      $\rho = \dfrac{800}{0.32} = \SI{2500}{\kilogram\per\metre\cubed}$

\item $V = \tfrac43 \pi r^{3} = \tfrac43 \pi (0.005)^{3} = \SI{5.24e-7}{\metre\cubed}$\\
      $m = \rho V = 7860 \times 5.24\times10^{-7} = \SI{0.00412}{\kilogram}$

\item (Triangular prism)\\
      $V = \dfrac{m}{\rho} = \dfrac{288000}{2400} = \SI{120}{\metre\cubed}$\\
      Cross‐sectional area $A = \dfrac{V}{l} = \dfrac{120}{5} = \SI{24}{\metre\squared}$\\
      $A = \tfrac12 b h \Rightarrow h = \dfrac{2A}{b} = \dfrac{2 \times 24}{6} = \SI{8}{\metre}$

\item (Sphere) $V = \dfrac{m}{\rho} = \dfrac{3966}{1850} = \SI{2.14}{\metre\cubed}$\\
      $r = \left(\dfrac{3V}{4\pi}\right)^{1/3} \approx \SI{0.800}{\metre}$

\item (Geode crust) \\
      $r_\text{outer} = \SI{0.04}{\metre},\; r_\text{inner} = 0.04 - 0.005 = \SI{0.035}{\metre}$\\
      $V_\text{crust} = \tfrac43\pi (r_\text{outer}^{3} - r_\text{inner}^{3})
        = \SI{8.86e-5}{\metre\cubed}$\\
      $m = \rho V = 1760 \times 8.86\times10^{-5} = \SI{0.156}{\kilogram}$
\end{enumerate}

\end{document}
